%% Пример резюме на русском - [YOUR_NAME]
%% Полный обзор компетенций, опыта и достижений.
%% Мастер-резюме, из которого /apply делает версии под конкретные вакансии.
%%
%% SETUP: запусти /setup, чтобы заполнить его своими данными.
%% Сборка: lualatex main_example_ru.tex
%% Кириллица: шрифт moderncv по умолчанию (Latin Modern) русских букв не содержит,
%% поэтому ниже подключён CMU (Computer Modern Unicode) — тот же рисунок с кириллицей
%% и знаком ₽. pdflatex для русского резюме не поддерживается.

\documentclass[11pt,a4paper,sans]{moderncv}
\moderncvstyle{banking}
\moderncvcolor{blue}

% Имя и заголовки разделов — цветом moderncv (color1); см. комментарий в main_example.tex.
\renewcommand*{\namefont}{\fontsize{34}{36}\bfseries\upshape}
\colorlet{firstnamecolor}{color1}
\colorlet{lastnamecolor}{color1}
\colorlet{namecolor}{color1}
\renewcommand*{\sectionstyle}[1]{{\sectionfont\color{color1}#1}}

\usepackage[utf8]{inputenc}
% Шрифты с кириллицей (lualatex + fontspec). Файлы cmun*.otf — пакет cm-unicode:
% есть в полном TeX Live и MiKTeX; в TinyTeX/BasicTeX: tlmgr install cm-unicode.
\usepackage{fontspec}
\setmainfont{cmunrm.otf}[BoldFont=cmunbx.otf, ItalicFont=cmunti.otf, BoldItalicFont=cmunbi.otf]
\setsansfont{cmunss.otf}[BoldFont=cmunsx.otf, ItalicFont=cmunsi.otf, BoldItalicFont=cmunso.otf]
% В CMU нет знака ₽ — берём его из Raleway, который лежит в репозитории
% (путь от cv/: резюме всегда собирается командой `cd cv && lualatex ...`).
\usepackage{newunicodechar}
\newfontfamily\rublefont{Raleway-Medium.otf}[Path=../cover_letters/OpenFonts/fonts/raleway/]
\newunicodechar{₽}{{\rublefont ₽}}
% Русские переносы; english — для фрагментов латиницей (названия технологий, ссылки).
\usepackage[english,russian]{babel}
\AtEndPreamble{\hypersetup{
    colorlinks=true,
    linkcolor=blue,
    filecolor=magenta,
    urlcolor=blue,
    pdftitle={[YOUR_NAME] - резюме},
    pdfpagemode=UseNone,
}}
\usepackage[scale=0.80]{geometry}
\usepackage{import}

% личные данные
\name{[Имя]}{[Фамилия]}
% Если адрес указывать не нужно, УДАЛИ эту строку целиком: \address{}{}{}
% падает с «There's no line here to end» во всех версиях moderncv.
\address{[Город, Россия]}{}{}
\phone[mobile]{[+7 XXX XXX-XX-XX]}
\email{[your.email@example.com]}
\extrainfo{\href{[https://t.me/your-username]}{Telegram}, \href{[https://github.com/your-username]}{GitHub}, \href{[https://career.habr.com/your-username]}{Хабр Карьера}}

\begin{document}

\makecvtitle

% ============================================================
%     О СЕБЕ
% ============================================================

\vspace{6pt}
\small{\textbf{Желаемая должность:} [Python-разработчик (Middle+)]. \textbf{Ожидания:} [от 250\,000~₽ на руки], [удалённо или гибрид, Москва].\\[3pt]
[3--5 строк о себе под конкретную роль: чем вы сильны именно для неё. Пример: «Backend-разработчик с 5 годами опыта в финтехе. Проектирую и поддерживаю высоконагруженные сервисы на Python и Go, довожу задачи от постановки до продакшена, выстраиваю мониторинг и CI/CD».]}

% ============================================================
%     КЛЮЧЕВЫЕ НАВЫКИ
% ============================================================

\section{Ключевые навыки}
\vspace{1pt}
\begin{itemize}

\item \textbf{[Категория навыков 1]}: [конкретные технологии, фреймворки, инструменты]

\item \textbf{[Категория навыков 2]}: [конкретные технологии, фреймворки, инструменты]

\item \textbf{[Категория навыков 3]}: [конкретные технологии, фреймворки, инструменты]

\item \textbf{[Категория навыков 4]}: [предметная экспертиза, методы, подходы]

\item \textbf{[Категория навыков 5]}: [платформы, сервисы, ПО]

\end{itemize}

% ============================================================
%     ОПЫТ РАБОТЫ
% ============================================================

% Даты — один ASCII-дефис ("2022-н.в."), не "--": двойной дефис становится
% коротким тире, и многие ATS перестают видеть диапазон (см. 05-cv-templates.md).
\section{Опыт работы}
\vspace{3pt}
\begin{itemize}

% --- Текущее место ---
\item{\cventry{[2022-н.в.]}{[Должность]}{[Компания]}{[Город]}{}{\vspace{1pt}
\begin{itemize}
    \item {[Достижение или зона ответственности 1 — конкретно, с цифрами, где возможно]}
    \item {[Достижение или зона ответственности 2]}
    \item {[Достижение или зона ответственности 3]}
    \item {[Достижение или зона ответственности 4]}
\end{itemize}}}

\vspace{3pt}

% --- Предыдущее место ---
\item{\cventry{[2019-2022]}{[Должность]}{[Компания]}{[Город]}{}{\vspace{1pt}
\begin{itemize}
    \item {[Достижение или зона ответственности 1]}
    \item {[Достижение или зона ответственности 2]}
    \item {[Достижение или зона ответственности 3]}
\end{itemize}}}

\vspace{3pt}

% --- Более раннее место ---
\item{\cventry{[2017-2019]}{[Должность]}{[Компания]}{[Город]}{}{\vspace{1pt}
\begin{itemize}
    \item {[Достижение или зона ответственности 1]}
    \item {[Достижение или зона ответственности 2]}
\end{itemize}}}

\end{itemize}

% ============================================================
%     ОБРАЗОВАНИЕ
% ============================================================

\section{Образование}
\vspace{1pt}
\begin{itemize}

\item{\cventry{[2015-2017]}{[Магистратура], [направление]}{[Вуз]}{[Город]}{}{\vspace{1pt}
Диссертация: «[Тема].» [Коротко о фокусе исследования.]
}}

\vspace{3pt}

\item{\cventry{[2011-2015]}{[Бакалавриат], [направление]}{[Вуз]}{[Город]}{}{\vspace{1pt}
[Коротко: ключевые дисциплины или проекты.]
}}

\end{itemize}

% ============================================================
%     ЯЗЫКИ
% ============================================================

\section{Языки}
\vspace{1pt}
\begin{itemize}
\item {Русский — родной, английский — [B2 (Upper-Intermediate)], [третий язык] — [уровень].}
\end{itemize}

% ============================================================
%     ПУБЛИКАЦИИ (необязательно)
% ============================================================

\section{Публикации}
\vspace{1pt}
\begin{itemize}
\item {[Автор(ы)] ([Год]). [Название]. [Журнал/конференция]. \href{[DOI_URL]}{DOI}}
\end{itemize}

% ============================================================
%     НАГРАДЫ (необязательно)
% ============================================================

\section{Награды и достижения}
\vspace{1pt}
\begin{itemize}
\item{\textbf{[Название награды]} -- [Мероприятие/организация] ([Год]).}
\end{itemize}

% ============================================================
%     РЕКОМЕНДАЦИИ
% ============================================================

\section{Рекомендации}
\vspace{1pt}
\begin{itemize}
\item{Рекомендации предоставлю по запросу.}
\end{itemize}

\end{document}
